\documentclass[11pt]{article}
\usepackage[a4paper,margin=25mm]{geometry}
\usepackage{amsmath,amssymb,amsthm,mathtools}
\usepackage{algorithm}
\usepackage{algpseudocode}
\usepackage{hyperref}
\usepackage{booktabs}

\newtheorem{definition}{Definition}
\newtheorem{assumption}{Assumption}
\newtheorem{lemma}{Lemma}
\newtheorem{theorem}{Theorem}
\newtheorem{proposition}{Proposition}
\newtheorem{corollary}{Corollary}

\title{QWeave Mathematical Foundations and Proof Obligations}
\author{Shrivardhini N \and Atharva Sheersh Pandey \and Diptesh Das \and Haridasu Sreedhar}
\date{}

\begin{document}
\maketitle

\section{Scope of the formal claims}
These results establish mapping preservation, hardware legality, semantic equivalence up to a declared wire permutation, schedule feasibility, and termination when the deterministic fallback assumptions hold. They do not establish global depth optimality or general convergence of PPO. Depth improvement is an empirical claim evaluated against matched baselines.

\section{Model}
\begin{definition}[Logical circuit]
Let $C=(V_C,E_C)$ be a finite dependency DAG. Each operation $v\in V_C$ has logical operands $Q_L(v)$ and duration $\tau(v)>0$ after target decomposition. The logical two-qubit interaction graph is $G_C=(L,E_I)$ with weights $w_{ij}\ge 0$.
\end{definition}

\begin{definition}[Hardware target]
The target is a directed graph $G_H=(P,E_H)$. An edge $(p,q)\in E_H$ exists only if the required native two-qubit instruction is supported from $p$ to $q$. Instruction duration and optional error are edge attributes.
\end{definition}

\begin{definition}[Mapping]
At routing step $t$, $\pi_t:L_{\mathrm{live}}\hookrightarrow P$ is injective. A logical two-qubit gate $(i,j)$ is executable when the required directed instruction is supported on $(\pi_t(i),\pi_t(j))$ after any validated direction correction.
\end{definition}

The initial mapping objective is
\begin{equation}
C_{\mathrm{map}}(\pi)=\sum_{(i,j)\in E_I}w_{ij}d_H\!\left(\pi(i),\pi(j)\right).
\end{equation}

\section{Assumptions}
\begin{assumption}[Finite input]
The circuit contains finitely many operations and the routing look-ahead is finite.
\end{assumption}
\begin{assumption}[Usable target component]
The chosen hardware component is connected when direction is ignored, contains at least $|L|$ usable physical qubits, and every fallback SWAP has a valid native decomposition.
\end{assumption}
\begin{assumption}[Guarded execution]
Every policy action is checked by a deterministic target guard before mutation of the circuit or mapping.
\end{assumption}
\begin{assumption}[Bounded non-progress]
After at most $K$ non-progress actions or a repeated-state threshold, control passes to the deterministic fallback.
\end{assumption}

\section{Initial mapping}
\begin{proposition}[Monotone pair-swap refinement]
If the local-refinement algorithm accepts a mapping $\pi'$ only when $C_{\mathrm{map}}(\pi')<C_{\mathrm{map}}(\pi)$, then it terminates after finitely many accepted moves at a pair-swap local optimum.
\end{proposition}
\begin{proof}
There are finitely many injections from the finite logical set into the finite physical set. Strict decrease prevents revisiting any accepted mapping. Hence only finitely many moves are accepted. At termination no pair swap in the selected neighbourhood has negative objective difference, which is precisely pair-swap local optimality. This does not imply global optimality.
\end{proof}

\section{Masked routing policy}
For a binary legality mask $m(s,a)$ and logit $z_\theta(s,a)$, define
\begin{equation}
\widetilde\pi_\theta(a\mid s)=
\frac{m(s,a)e^{z_\theta(s,a)}}{\sum_{a'\in\mathcal A}m(s,a')e^{z_\theta(s,a')}}.
\end{equation}

\begin{lemma}[Illegal-action exclusion]
If at least one legal action exists, then $m(s,a)=0$ implies $\widetilde\pi_\theta(a\mid s)=0$.
\end{lemma}
\begin{proof}
The numerator is zero for every masked action. The denominator is positive because at least one action has mask value one and a finite exponential logit.
\end{proof}

\begin{lemma}[Mapping preservation under SWAP]
If $\pi_t$ is injective and a SWAP exchanges the occupants of two distinct physical vertices, then $\pi_{t+1}$ is injective.
\end{lemma}
\begin{proof}
The physical exchange is a transposition $\sigma:P\to P$, hence a bijection. The updated map is $\pi_{t+1}=\sigma\circ\pi_t$. The composition of an injection and a bijection is injective.
\end{proof}

\section{Hardware legality and semantics}
\begin{theorem}[Hardware legality]
Every two-qubit operation emitted by QWeave is supported by the target on the scheduled physical pair.
\end{theorem}
\begin{proof}
The only routing mutations are guarded SWAPs and guarded execution of frontier gates. A SWAP is emitted only on a coupling edge with a validated native decomposition. A logical two-qubit gate is emitted only when its mapped ordered pair supports the required instruction or a validated direction-correction sequence. Scheduling changes start times but not operands. Therefore every emitted two-qubit operation is target legal.
\end{proof}

\begin{theorem}[Semantic equivalence up to wire permutation]
Let $U_C$ be the logical circuit unitary, $P_{\pi_0}$ the initial layout permutation, and $P_{\pi_T}$ the final layout permutation. Ignoring measurements until their preserved terminal positions, the routed circuit satisfies
\begin{equation}
U_{\mathrm{out}}=P_{\pi_T}U_CP_{\pi_0}^{\dagger}.
\end{equation}
\end{theorem}
\begin{proof}
Proceed by induction over emitted operations. Initially the physical interpretation is $\pi_0$. For an emitted logical gate, the guard ensures it acts on the physical images of the same logical operands, so the logical action is preserved. For an inserted SWAP, the quantum states at its two physical locations are exchanged and the classical mapping is updated by the identical transposition. Thus the interpretation invariant is preserved after every step. Composition of all steps yields the logical unitary conjugated by the initial and final wire permutations.
\end{proof}

\section{Termination}
\begin{theorem}[Termination with shortest-path fallback]
Under the stated assumptions, QWeave terminates after emitting a finite routed circuit.
\end{theorem}
\begin{proof}
Policy non-progress is bounded by $K$, after which fallback is invoked. Select a ready two-qubit gate with mapped operands $p$ and $q$. Because the usable target component is connected, a finite path from $p$ to $q$ exists. The fallback applies a legal SWAP that moves one operand along a shortest path, decreasing the remaining path length by one. After finitely many such swaps the operands are adjacent and the gate executes, decreasing the number of unexecuted logical gates. Since the input circuit is finite, repeated retirement of a gate terminates the procedure.
\end{proof}

\section{Scheduling}
For ready operation $v$, define the earliest start time
\begin{equation}
s(v)=\max\left\{\max_{u\prec v}(s(u)+\tau(u)),\ \max_{q\in Q(v)}A(q)\right\},
\end{equation}
where $A(q)$ is the availability time of physical qubit $q$.

\begin{theorem}[Schedule feasibility]
The resulting schedule respects all routed-DAG precedences and contains no overlap between operations sharing a physical qubit.
\end{theorem}
\begin{proof}
Only ready operations are scheduled, and the first maximum ensures every predecessor finishes before $v$ starts. The second maximum ensures every operand is available. After placement, setting $A(q)\leftarrow s(v)+\tau(v)$ for every operand prevents any later operation on those qubits from overlapping $v$.
\end{proof}

\begin{proposition}[Lower bound]
Any feasible schedule has makespan
\begin{equation}
D^\star\ge \max\left\{\operatorname{CP}(C),\ \max_{q\in P}\sum_{v:q\in Q(v)}\tau(v)\right\}.
\end{equation}
\end{proposition}
\begin{proof}
Every precedence path must execute sequentially, giving the critical-path bound. Every operation using a fixed qubit is mutually exclusive on that resource, giving the load bound. Their maximum is therefore a valid lower bound.
\end{proof}

\section{Reward shaping and PPO boundary}
Use
\begin{equation}
r_t=-\Delta D_t-\lambda_S\Delta S_t-\lambda_E\Delta E_t+\gamma\Phi(s_{t+1})-\Phi(s_t).
\end{equation}
The final term is potential-based shaping. Under the usual discounted-MDP assumptions, the standard policy-invariance result of Ng, Harada, and Russell applies. PPO maximises a clipped empirical surrogate and is not covered by a general global-convergence or global-depth-optimality theorem. Accordingly, QWeave reports learned performance through held-out experiments.

\section{Implementation-linked proof obligations}
\begin{center}
\begin{tabular}{p{0.28\linewidth}p{0.31\linewidth}p{0.31\linewidth}}
\toprule
Claim & Runtime assertion & Required test\\
\midrule
Injective mapping & no duplicate physical values & property test over random SWAP sequences\\
Legal execution & target instruction lookup succeeds & adversarial directed-coupling cases\\
Fallback progress & shortest-path distance decreases & paths, grids, irregular connected graphs\\
Schedule feasibility & precedence and interval audit & random DAG/resource conflicts\\
Semantic equivalence & final-layout-aware operator check & random small circuits and metamorphic rewrites\\
\bottomrule
\end{tabular}
\end{center}

\section{Claim discipline}
The paper may state that QWeave is correct and terminating only under the declared assumptions and only when the corresponding executable checks pass. It may state that QWeave reduces depth only with a quantitative result, uncertainty interval, matched baseline, and disclosed failure count. It must not infer global optimality from PPO training or from a finite benchmark.

\begin{thebibliography}{9}
\bibitem{li2019sabre} G. Li, Y. Ding, and Y. Xie, ``Tackling the Qubit Mapping Problem for NISQ-Era Quantum Devices,'' ASPLOS, 2019.
\bibitem{schulman2017ppo} J. Schulman et al., ``Proximal Policy Optimization Algorithms,'' arXiv:1707.06347, 2017.
\bibitem{gilmer2017mpnn} J. Gilmer et al., ``Neural Message Passing for Quantum Chemistry,'' ICML, 2017.
\bibitem{ng1999shaping} A. Y. Ng, D. Harada, and S. Russell, ``Policy Invariance under Reward Transformations,'' ICML, 1999.
\bibitem{quetschlich2023mqt} N. Quetschlich, L. Burgholzer, and R. Wille, ``MQT Bench,'' Quantum 7, 1062, 2023.
\end{thebibliography}
\end{document}
